\kbvsechead{ಅಕ್ಷರ ಗಣಗಳು :}

\indent ಅಕ್ಷರ ಗಣಗಳಲ್ಲಿ ಅಕ್ಷರಗಳ ಸಂಖ್ಯೆ ಗಣನೀಯವೇ ಹೊರತು. ಮೂರು ಅಕ್ಷರಗಳ ಗುಂಪು ಒಂದು ಗಣವೆನಿಸುವುದು. ಆದುದರಿಂದ ಪದ್ಯದ ಪ್ರತಿ ಸಾಲಿನಲ್ಲಿರುವ ಅಕ್ಷರಗಳನ್ನು ಮೂರು ಮೂರು ಅಕ್ಷರಗಳಂತೆ ಲೆಕ್ಕಮಾಡಿ ಗಣಗಳನ್ನು ವಿಂಗಡಿಸಲಾಗುವುದು. ಈ ಬಗೆಯ ವಿಂಗಡನೆಯು “ವೃತ್ತ ಜಾತಿಯ ಛಂದಸ್ಸಿಗೆ” ಸೀಮಿತವಾಗಿರುವುದು. ವೃತ್ತ ಜಾತಿಯ ಛಂದಸ್ಸಿನಲ್ಲಿ (೧) ಉತ್ಪಲಮಾಲಾ ವೃತ್ತ, (೨) ಚಂಪಕಮಾಲಾ ವೃತ್ತ, (೩) ಶಾರ್ದೂಲ ವಿಕ್ರೀಡಿತ ವೃತ್ತ, (೪) ಮತ್ತೇಭ ವಿಕ್ರೀಡಿತ ವೃತ್ತ, (೫) ಸ್ರಗ್ಧರಾ ವೃತ್ತ, (೬) ಮಹಾಸ್ರಗ್ಧರಾ ವೃತ್ತ ಎಂಬ ಆರು ಜಾತಿಯ ವೃತ್ತಗಳು ಮುಖ್ಯವಾದವುಗಳು.\par

\indent ಅಕ್ಷರ-ಗಣಗಳಲ್ಲಿ ಎಂಟು ವಿಧ: ಯಗಣ, ಮಗಣ, ತಗಣ, ರಗಣ, ಜಗಣ, ಭಗಣ, ನಗಣ ಮತ್ತು ಸಗಣ-ಇವೇ ಆ ಎಂಟು ಗಣಗಳು. “ಯಮಾತಾರಾಜಭಾನಸಲಗಂ” ಎಂಬ ಅಕ್ಷರ ಸೂತ್ರದ ಸಹಾಯದಿಂದ ಮೇಲೆ ತಿಳಿಸಿದ ಗಣಗಳ ಗುರು-ಲಘು ವಿನ್ಯಾಸವನ್ನು ತಿಳಿಯಬಹುದು. ಸೂತ್ರದಲ್ಲಿನ ಯಕಾರದಿಂದ ಪ್ರಾರಂಭಿಸಿ ಪ್ರತಿ ಮೂರು ಕ್ರಮಗತ ಅಕ್ಷರಗಳಿಗೆ ಪ್ರಸ್ತಾರ ಹಾಕಿದರೆ ಯಗಣದಿಂದ ಸಗಣದವರೆಗಿನ ಅಕ್ಷರ ಗಣಗಳ ವಿನ್ಯಾಸ ದೊರಕುತ್ತದೆ.\par

\kbvsechead{ಗಣಗಳ ವಿನ್ಯಾಸ}

\vspace{3pt plus 2pt}\noindent\hspace*{1.2em}\kbvlabel{ಉದಾ}\par\vspace{1pt}

\vspace{5pt plus 4pt minus 1pt}
\noindent{\renewcommand{\arraystretch}{1.24}\setlength{\tabcolsep}{4pt}\rowcolors{1}{kbvrow}{white}
\begin{tabularx}{\textwidth}{X X X X X}
ಗಣ & ಸೂತ್ರ & ವಿನ್ಯಾಸ & ಲಕ್ಷಣ & ಉದಾಹರಣೆ \\
೧. ಯಗಣ & ಯಮಾತಾ & $\mbox{\laghu}$ \guru{} \guru{} & (ಆದಿಲಘು) & : ಯಥಾರ್ಥಂ \\
೨. ಮಗಣ & ಮಾತಾರಾ & \guru{} \guru{} \guru{} & (ಸರ್ವಗುರು) & : ಮಸ್ತಿಷ್ಕಂ \\
೩. ತಗಣ & ತಾರಾಜ & \guru{} \guru{} $\mbox{\laghu}$ & (ಅಂತ್ಯಲಘು) & : ತಣ್ಣೀರು \\
೪. ರಗಣ & ರಾಜಭಾ & \guru{} $\mbox{\laghu}$ \guru{} & (ಮಧ್ಯಲಘು) & : ರಕ್ಕಸೀ \\
೫. ಜಗಣ & ಜಭಾನ & $\mbox{\laghu}$ \guru{} $\mbox{\laghu}$ & (ಮಧ್ಯಗುರು) & : ಜಟಾಯು \\
೬. ಭಗಣ & ಭಾನಸ & \guru{} $\mbox{\laghu}$ $\mbox{\laghu}$ & (ಆದಿಗುರು) & : ಭವ್ಯತೆ \\
೭. ನಗಣ & ನಸಲ & $\mbox{\laghu}$ $\mbox{\laghu}$ $\mbox{\laghu}$ & (ಸರ್ವಲಘು) & : ನಕುಲ \\
೮. ಸಗಣ & ಸಲಗಂ & $\mbox{\laghu}$ $\mbox{\laghu}$ \guru{} & (ಅಂತ್ಯಗುರು) & : ಸರ್ವಥಾ \\
\end{tabularx}}
\vspace{5pt plus 4pt minus 1pt}
